\documentclass[3p,twocolumn]{elsarticle}

% ==============================================================================
% Packages
% ==============================================================================
\usepackage[utf8]{inputenc}
\usepackage{amsmath,amssymb,amsfonts}
\usepackage{graphicx}
\usepackage{booktabs}
\usepackage{multirow}
\usepackage{array}
\usepackage{url}
\usepackage{xcolor}
\usepackage{algorithm}
\usepackage{algorithmic}

\newcommand{\etal}{\textit{et al.}}
\newcommand{\eg}{\textit{e.g.}, }
\newcommand{\ie}{\textit{i.e.}, }

\setlength{\textfloatsep}{9pt plus 2pt minus 2pt}
\setlength{\floatsep}{8pt plus 2pt minus 2pt}
\setlength{\intextsep}{8pt plus 2pt minus 2pt}

\usepackage[colorlinks=true,
            linkcolor=blue!70!black,
            citecolor=blue!70!black,
            urlcolor=blue!70!black]{hyperref}

\graphicspath{{final_deliverables/figures/}{../results/training_multiclass/}{../results/vehicle_dynamics/}{../results/safety_analysis/}{../results/ncap_testing/}}

\journal{Transportation Research Part C: Emerging Technologies}

% ==============================================================================
% Document
% ==============================================================================
\begin{document}

\begin{frontmatter}

\title{A Physics-Aware Vision-Based Automatic Emergency Braking System for Vulnerable Road User Protection: Multi-Class Perception, Adaptive Time-to-Collision Control, Brake Thermodynamics, and Biomechanical Injury Mitigation}

\author[iitr]{Atharv Priyadarshi}
\author[iitr]{Agrawal Chehek Gopal}
\author[iitr]{Debangan Sarkar}
\author[iitr]{Arnav Garg}
\author[iitr]{Indra Vir Singh\corref{cor1}}
\ead{ivsingh@me.iitr.ac.in}

\cortext[cor1]{Corresponding author}
\address[iitr]{Department of Mechanical and Industrial Engineering, Indian Institute of Technology Roorkee, Roorkee 247667, Uttarakhand, India}

\begin{abstract}
Vulnerable Road Users (VRUs)---including pedestrians, cyclists, and motorcyclists---constitute over 26\% of global traffic fatalities, with developing countries such as India reporting pedestrian death shares exceeding 15\% of total road casualties. Automatic Emergency Braking (AEB) provides a vital active safety intervention; however, contemporary production systems overwhelmingly rely on simplified point-mass kinematic formulations ($d = v^2 / 2\mu g$) that treat tyre grip as a linear constant, disregard brake rotor thermal fade, and assume invariant driver situational awareness. In real-world emergency stops, these assumptions systematically underestimate stopping distances by 15\% to 30\%, precipitating delayed or inadequate deceleration commands. This investigation presents an end-to-end, physics-aware, monocular vision-based AEB pipeline that directly couples real-time environmental perception with non-linear vehicle dynamics, closed-loop brake thermodynamics, in-cabin cognitive monitoring, and biomechanical injury assessment. The perception layer employs a fine-tuned nano-scale YOLOv8 model on a stratified 6,000-image split of the BDD100K driving dataset, achieving an overall mAP@0.5 of 40.0\%, an mAP@0.5:0.95 of 21.3\%, a precision of 53.6\%, and a recall of 37.9\% across four target traffic classes, running at 15--20 frames per second on an NVIDIA Tesla T4. Persistent identity association is executed via Kalman-filtered SORT tracking, and monocular metric distance is resolved using calibrated pinhole geometry and Bird's-Eye View (BEV) mapping. Longitudinal stopping distances are modeled using the Pacejka '92 Magic Formula tyre model integrated at a 1\,ms timestep, incorporating dynamic axle load transfers and closed-loop anti-lock braking system (ABS) modulation across dry ($\mu=0.85$), wet ($\mu=0.55$), gravel ($\mu=0.45$), snow ($\mu=0.25$), and ice ($\mu=0.10$) surfaces. A thermodynamic brake rotor model continuously tracks friction degradation using kinetic energy conservation and Newton's convective cooling law, penalizing the Time-to-Collision (TTC) activation threshold by up to $+3.0$\,s when rotors experience severe thermal fade ($T \ge 650^\circ$C). In-cabin driver fatigue is tracked via 68-point facial landmark Eye Aspect Ratio (EAR) and 3D Perspective-n-Point head pose tracking, dynamically scaling reaction delays from 0.7\,s (alert) to 1.6\,s (drowsy) and applying an immediate throttle override under microsleep conditions. Regulatory compliance is verified through the standardized Euro NCAP Car-to-Pedestrian protocol across four scenarios (CPFA, CPNA, CPNC, CPLA) and nine approach speeds (20--60\,km/h), yielding an overall performance score of 97\% on dry asphalt (34 avoided, 1 partial, 1 collision) and 92\% on wet asphalt. Biomechanical head impact kinematics are evaluated through the Head Injury Criterion (HIC) and the Abbreviated Injury Scale (AIS), demonstrating that partial speed attenuation from 50 to 30\,km/h mitigates injury from fatal/critical levels down to moderate head trauma. The pipeline further integrates heuristic and neural traffic sign recognition for regulatory speed zone compliance, Grad-CAM visual explainability maps for perception accountability, and a circular black-box incident telemetry recorder. This paper provides an open, fully verifiable benchmark establishing that physics-grounded perception-control coupling is both technically feasible and essential for resilient autonomous vehicle safety.
\end{abstract}

\begin{keyword}
Automatic emergency braking (AEB) \sep Vulnerable road users (VRUs) \sep YOLOv8 object detection \sep Pacejka tyre model \sep Brake thermodynamics \sep Driver monitoring systems \sep Euro NCAP protocol \sep Head Injury Criterion (HIC) \sep Visual explainability
\end{keyword}

\end{frontmatter}

% ==============================================================================
\section{Introduction}
\label{sec:intro}
% ==============================================================================

Preventable road traffic trauma represents an urgent public health crisis. According to the World Health Organization's Global Status Report on Road Safety, approximately 1.19~million individuals lose their lives annually on global roadways~\cite{who2023}. Vulnerable Road Users (VRUs)---specifically pedestrians, pedal cyclists, and operators of two- and three-wheelers---account for disproportionate casualty shares, exceeding 26\% of global mortality. In rapidly motorizing developing nations with heterogeneous traffic mixtures, the vulnerability is even more pronounced: official data from India's Ministry of Road Transport and Highways (MoRTH) recorded 168,491 fatalities in 2022, with pedestrians constituting 19.5\% (32,825 deaths) and two-wheeler riders comprising 44.5\%~\cite{morth2023}. Unlike vehicle occupants protected by crumple zones, pre-tensioned seatbelts, and supplemental inflatable restraints, pedestrians absorb collision energy directly into soft tissue and the musculoskeletal frame. Consequently, active collision avoidance technology capable of detecting human obstacles and autonomously arresting vehicle motion represents the single most effective technological countermeasure.

Modern Advanced Driver Assistance Systems (ADAS) and Level~2+ autonomous vehicles rely heavily on Autonomous Emergency Braking (AEB) architectures mandated by international regulatory bodies, including UNECE Regulation No.~152~\cite{unece2020} and Euro NCAP assessment protocols~\cite{euroncap2023}. Despite widespread deployment, existing commercial and academic AEB algorithms exhibit fundamental systemic weaknesses originating from their reliance on point-mass kinematic formulations:
\begin{equation}
  d_{\text{stop}} \;=\; \frac{v^2}{2\,\mu\,g},
  \label{eq:kinematic}
\end{equation}
where $v$ denotes forward vehicle velocity, $\mu$ is an assumed static tyre-road friction coefficient, and $g = 9.81\,\text{m/s}^2$ is the acceleration due to gravity. 

Although analytically tractable, Eq.~\eqref{eq:kinematic} embeds four severe engineering simplifications that fail catastrophically in safety-critical edge cases:
\begin{enumerate}
  \item \textbf{Tyre-Force Non-Linearity and Slip Dynamics}: The friction coefficient $\mu$ is treated as an instantaneous, load-invariant constant. In physical vehicle tyres, longitudinal shear force exhibits non-linear saturation across varying longitudinal slip ratios ($\kappa$), governed by semi-empirical contact mechanics such as the Pacejka Magic Formula~\cite{pacejka1992,pacejka2005}. Peak retarding traction occurs only within an optimal slip window ($\kappa \in [0.10, 0.15]$). Under panic pedal application without dynamic ABS modeling, tyres lock up, sliding friction degrades, and stopping distance extends substantially.
  \item \textbf{Dynamic Axle Load Redistribution}: Severe deceleration induces forward inertial pitch moments, transferring vertical normal load from the rear axle to the front axle ($F_{z,\text{front}} > F_{z,\text{rear}}$). Because tyre force generation is load-dependent and exhibits non-linear vertical load sensitivity, ignoring weight transfer skews the optimal brake force allocation between front and rear calipers~\cite{rajamani2011,gillespie1992}.
  \item \textbf{Brake Rotor Thermal Degradation (Fade)}: Under repeated hard decelerations or prolonged downhill driving, kinetic energy transformed into thermal energy elevates cast-iron rotor surfaces to temperatures exceeding $400^\circ\text{C}$ to $650^\circ\text{C}$. Automotive friction pairs undergo thermal fade, wherein resin binders vaporize and pad-rotor friction coefficients collapse by 30\% to 70\%~\cite{day2014,tirovic2007,limpert2011}. Standard ADAS controllers assume constant cold-friction performance ($\mu_{\text{pad}} \approx 0.40$), executing AEB deceleration profiles that fail to arrest the vehicle before impact.
  \item \textbf{Human Cognitive State Decoupling}: Standard AEB control logic utilizes fixed Time-to-Collision (TTC) triggering thresholds (typically 1.2 to 1.5\,s). However, human drivers exhibit widely varying reaction delays depending on drowsiness, ocular closure, and gaze distraction~\cite{soukupova2016,dua2013}. A static intervention threshold that affords an alert driver sufficient agency to initiate manual avoidance becomes lethal when the driver is fatigued or experiencing microsleep.
\end{enumerate}

Physical stopping distance experiments and rigorous multi-body vehicle simulations confirm that these compounded simplifications cause standard kinematic estimators to underestimate true stopping distances by 15\% to 30\% on dry asphalt, and by more than 40\% on contaminated, wet, or icy pavements~\cite{iso21994,seiniger2012}.

To resolve these deficiencies, this research develops, verifies, and evaluates an open-source, physics-aware, vision-based AEB pipeline. The framework unifies environmental computer vision, vehicle contact mechanics, brake thermodynamics, driver cabin monitoring, regulatory validation, and crash biomechanics into a tightly coupled real-time architecture.

\subsection{Principal Contributions}
The major contributions of this paper are:
\begin{itemize}
  \item \textbf{Empirically Grounded Perception Benchmark}: We present a complete, uninflated performance audit of a nano-scale deep neural network (YOLOv8n) fine-tuned on a 6,000-image stratified split of BDD100K across four VRU-critical classes. We candidly document actual convergence metrics (mAP@0.5 of 40.0\%, mAP@0.5:0.95 of 21.3\%, precision of 53.6\%, and recall of 37.9\%), analyzing the fundamental trade-offs between embedded edge inference latency (15--20\,fps on Tesla T4) and small-object bounding-box recall.
  \item \textbf{Pacejka Longitudinal Dynamics Simulator}: We integrate a five-force dynamic braking simulator executing at a 1\,ms discrete timestep. Incorporating Pacejka '92 tyre saturation, dynamic front-to-rear axle load transfer, aerodynamic drag, rolling resistance, and a 50\,ms cyclic ABS pressure modulation loop, the model reproduces physical stopping distances across dry, wet, gravel, snow, and ice surfaces that strictly adhere to ISO~21994 reference standards.
  \item \textbf{Closed-Loop Brake Thermodynamics Model}: We introduce an energy-conservation heating and Newton-cooling formulation that tracks real-time rotor thermal accumulation and maps empirical friction-fade degradation into an adaptive AEB TTC penalty of up to $+3.0$\,s---a mechanical capability completely absent from existing commercial production ADAS.
  \item \textbf{Human-Aware Cognitive Intervention Controller}: We implement a driver cabin monitoring subsystem utilizing 68-point facial landmark Eye Aspect Ratio (EAR) and solvePnP head-pose tracking. The module maps driver alertness states (Alert, Tired, Distracted, Drowsy, Microsleep) to adaptive TTC buffer extensions (0.7 to 1.6\,s) and an emergency throttle override.
  \item \textbf{Multi-Physics ML Classifier and Hybrid Safety Arbitration}: We train an ensemble Random Forest classifier ($N_{\text{trees}}=250$) across 12,000 multi-physics scenario corridors ($[\text{TTC}, d, v_{\text{rel}}, \mu_{\text{weather}}, \text{EAR}, T_{\text{rotor}}]$), achieving 98.67\% validation agreement with ISO~22839 corridors and 97.44\% test accuracy, coupled with a deterministic physical arbitration envelope to guarantee fail-safe intervention.
  \item \textbf{Euro NCAP 36-Test Matrix Verification}: The complete pipeline is evaluated across Euro NCAP's four standardized pedestrian scenarios (CPFA-50, CPNA-25, CPNC-50, CPLA) across nine velocities (20--60\,km/h) on both dry and wet surfaces, demonstrating 97\% and 92\% overall regulatory compliance scores.
  \item \textbf{Biomechanical Injury Mitigation \& Systems Explainability}: We bridge control actuation to clinical trauma by modeling Head Injury Criterion (HIC) and Abbreviated Injury Scale (AIS) curves, demonstrating that pre-crash partial braking yields dramatic non-linear injury reductions. We augment the pipeline with regulatory traffic sign recognition, Grad-CAM perceptual visual saliency, and a circular black-box crash recorder.
\end{itemize}

% ==============================================================================
% Table 1: Comprehensive Literature Comparison
% ==============================================================================
\begin{table*}[!t]
\caption{Systematic Architectural Comparison of Representative AEB Research Frameworks and Commercial Standards}
\label{tab:related_work}
\centering
\small
\setlength{\tabcolsep}{4.5pt}
\resizebox{\textwidth}{!}{%
\begin{tabular}{lcccccccc}
\toprule
\textbf{Framework / Reference} & \textbf{Perception Sensor} & \textbf{Detector Backbone} & \textbf{Tyre Model} & \textbf{Brake Fade} & \textbf{Driver State} & \textbf{Injury HIC} & \textbf{NCAP Suite} & \textbf{Latency} \\
\midrule
Seiniger \etal~\cite{seiniger2012}       & Radar / Lidar & Geometric Cluster & Kinematic ($v^2/2a$) & No & No & No & 4 Scenarios & N/A \\
Ger\'onimo \etal~\cite{geronimo2010}     & Stereo Vision & HOG + Linear SVM  & Kinematic & No & No & No & None & $>$100\,ms \\
L\"ubbe \etal~\cite{lubbe2018}           & Monocular CAM & Statistical Proxy & Kinematic & No & No & Empirical & Crash Data & Offline \\
Althoff \& Dolan~\cite{althoff2014}     & Multi-Sensor  & Occupancy Grid    & Linear Slip Model & No & No & No & Reachability & 25\,ms \\
Chugh \etal~\cite{chugh2021}             & Video Logs    & Faster R-CNN      & None (Perception) & No & No & Statistical & None & 85\,ms \\
Grover \etal~\cite{grover2022}           & Simulation    & FEA Biomechanics  & Contact Impact    & No & No & AIS / HIC & Component & Offline \\
Dua \etal~\cite{dua2013}                 & Cabin Camera  & Facial Geometry   & None (DMS Only)   & No & Alertness & No & None & 30\,ms \\
UNECE Reg 152~\cite{unece2020}           & OEM Discretion& Unspecified       & Prescribed Decel  & No & Invariant & Pass/Fail & Mandated & OEM \\
Euro NCAP Protocol~\cite{euroncap2023}   & Physical Test & Physical Dummies  & Real Vehicle      & Physical & Invariant & HIC Probe & 4 Scenarios & Real-time \\
\midrule
\textbf{Proposed System}                 & \textbf{Monocular CAM} & \textbf{YOLOv8n C2f} & \textbf{Pacejka '92 + ABS} & \textbf{Thermodynamic} & \textbf{EAR + Pose} & \textbf{Analytical HIC} & \textbf{36 Tests} & \textbf{52\,ms} \\
\bottomrule
\end{tabular}
}%
\end{table*}

% ==============================================================================
\section{Related Work and Comparative Taxonomy}
\label{sec:related}
% ==============================================================================

\subsection{Object Detection in Intelligent Transportation Systems}
Early pedestrian protection systems relied on hand-crafted spatial feature representations, prominently Histograms of Oriented Gradients (HOG) coupled with Support Vector Machines (SVMs)~\cite{dalal2005} and Deformable Part Models (DPMs)~\cite{felz2010}. While foundational, these classical detectors suffered from prohibitive false-positive rates in cluttered urban visual backgrounds and catastrophic latency degradation under multi-scale sliding-window searches~\cite{geronimo2010}. The introduction of deep convolutional architectures established transformative gains: two-stage frameworks such as Faster R-CNN~\cite{ren2017} introduced Region Proposal Networks (RPNs), achieving unprecedented bounding-box localization accuracy at the expense of high computational overhead ($\ge 80$\,ms per frame on desktop GPUs), rendering them impractical for low-power embedded automotive Electronic Control Units (ECUs).

Single-stage detectors, notably SSD~\cite{liu2016ssd} and the You Only Look Once (YOLO) lineage~\cite{redmon2018}, fundamentally unified bounding box regression and semantic classification into a single forward pass. Jocher \etal~\cite{jocher2023} introduced YOLOv8, incorporating an anchor-free split-head architecture and Cross-Stage Partial Bottleneck with two convolutions (C2f) modules. In this work, we deploy the nano variant (YOLOv8n), which possesses approximately 3.2\,million parameters and 8.7 GFLOPs at $640\times640$ image resolution. This architectural scale satisfies the strict sub-45\,ms perception budget demanded by closed-loop automotive actuation while maintaining multi-class detection across pedestrians, passenger cars, pedal cyclists, and motorcycles.

\subsection{Monocular Depth Estimation and Target Tracking}
Metric depth perception from a single camera presents an ill-posed inverse optics challenge. Self-supervised monocular depth estimation networks such as MiDaS~\cite{ranftl2022} have advanced dense scene depth reconstruction; nevertheless, they output relative, unscaled disparity maps that require auxiliary ground-truth sensors (e.g., LiDAR or radar) for metric conversion, and incur heavy neural inference latency ($>35$\,ms). For structured highway and urban driving scenes, the camera-ground pinhole projection model based on camera mounting elevation, pitch angle, and human anthropometric bounding box geometry yields robust metric depth estimates with sub-10\% longitudinal distance errors within the critical 30\,m AEB operating corridor.

For inter-frame temporal consistency, modern ADAS tracking relies on the tracking-by-detection paradigm. The Simple Online and Realtime Tracking (SORT) algorithm~\cite{bewley2016} couples a linear constant-velocity Kalman filter with the Hungarian bipartite assignment algorithm~\cite{kuhn1955} based on Intersection-over-Union (IoU) affinity. Although DeepSORT~\cite{wojke2017} and ByteTrack~\cite{zhang2022} introduce visual appearance feature embeddings and hierarchical low-threshold association, their neural re-identification backbones consume significant GPU cycles. For deterministic, sub-millisecond execution on edge automotive hardware, standard SORT provides the optimal computational trade-off.

\subsection{Longitudinal Vehicle Dynamics and Tyre Mechanics}
The mechanical efficacy of emergency braking is bounded by tyre-road interfacial friction. Classical automotive ADAS literature frequently abstracts braking to simple kinematic formulas ($d = v^2/2a$) or constant deceleration bounds ($a_{\text{max}} = \mu g$)~\cite{seiniger2012}. In dynamic tire mechanics, Hans Pacejka formulated the Magic Formula~\cite{pacejka1992,pacejka2005}, characterizing longitudinal retarding force ($F_x$) as an empirical trigonometric function of longitudinal slip ratio ($\kappa$), normal axle load ($F_z$), and road surface parameters ($B, C, D, E$). 

During severe deceleration, dynamic load transfer shifts vertical force onto the front steering axle while unloading the rear axle~\cite{rajamani2011}. If rear tyre forces exceed available traction, severe directional instability or spin-out occurs; conversely, front lock-up destroys steering control. Anti-lock Braking Systems (ABS) continually modulate brake caliper hydraulic pressure to constrain wheel slip within the stable positive-gradient zone ($\kappa \approx 0.10$--$0.15$). Despite the universality of ABS in production vehicles, prior vision-based AEB literature predominantly evaluates perception accuracy in isolation without executing a dynamic multi-force chassis simulation.

\subsection{Brake Thermodynamics and Friction Fade}
During an emergency stop from highway speeds, the vehicle's kinetic energy is converted almost entirely into thermal energy via frictional sliding between brake pads and cast-iron rotors~\cite{day2014}. The bulk temperature rise is governed by the energy balance equation:
\begin{equation}
  \Delta T \;=\; \frac{\Delta E_{\text{kin}}}{m_{\text{rotor}}\,c_p},
  \label{eq:heat_energy}
\end{equation}
where $m_{\text{rotor}}$ is the active thermal mass and $c_p$ is the specific heat capacity of gray cast iron ($\approx 450\,\text{J/(kg}\cdot\text{K)}$)~\cite{tirovic2007}. Under repetitive panic stops or high-speed mountain descents, surface temperatures routinely exceed $400^\circ\text{C}$ to $600^\circ\text{C}$. Automotive friction materials undergo thermal degradation: phenolic resins degrade, generating a microscopic gas layer between pad and disc that drastically lowers the friction coefficient $\mu_{\text{pad}}$~\cite{limpert2011}. Tirovic and Galindo-Lopez~\cite{tirovic2007} demonstrated that cooling occurs via velocity-dependent forced convection obeying Newton's law. To date, no commercial ADAS algorithm modulates its Time-to-Collision braking initiation threshold based on an online estimate of brake rotor temperature.

\subsection{Driver State Monitoring}
Real-world collision risk is a coupled function of environmental hazards and human cognitive capacity. Soukupov\'a and \v{C}ech~\cite{soukupova2016} established the Eye Aspect Ratio (EAR), derived from the Euclidean coordinates of facial landmarks surrounding the ocular palpebral fissure. Tracking EAR over contiguous frame windows enables reliable discrimination between voluntary rapid blinks ($\le 0.3$\,s), fatigue-induced extended closures ($> 0.5$\,s), and catastrophic microsleep episodes ($> 1.5$\,s). Integrating driver state estimation directly into AEB actuation logic transforms the vehicle from an oblivious mechanical actuator into a human-in-the-loop cognitive co-pilot.

\subsection{Injury Biomechanics and Regulatory Protocols}
Quantifying the clinical efficacy of an AEB system requires translating residual impact velocities into human tissue trauma metrics. The Head Injury Criterion (HIC), formulated by Versace~\cite{versace1971} and standardized in NHTSA FMVSS~208~\cite{eppinger1999} and Euro NCAP Pedestrian Testing Protocols~\cite{euroncap_ped_v9}, evaluates the temporal acceleration integral experienced by the cranial mass during impact:
\begin{equation}
  \text{HIC} \;=\; \max_{t_1, t_2} \left[ (t_2 - t_1) \left( \frac{1}{t_2 - t_1} \int_{t_1}^{t_2} a(t)\,\text{d}t \right)^{2.5} \right].
  \label{eq:hic_def}
\end{equation}
The Abbreviated Injury Scale (AIS)~\cite{aaam2016} classifies resulting trauma into six clinical tiers: AIS~1 (Minor), AIS~2 (Moderate), AIS~3 (Serious), AIS~4 (Severe), AIS~5 (Critical), and AIS~6 (Fatal). Table~\ref{tab:related_work} synthesizes the architectural divergence between prior literature and our unified framework.

% ==============================================================================
% Section 3: Architecture
% ==============================================================================
\section{End-to-End System Architecture}
\label{sec:arch}

\begin{figure*}[!t]
  \centering
  \includegraphics[width=0.88\textwidth]{fig_06_pipeline_flowchart}
  \caption{Architectural flowchart of the unified 14-step physics-aware AEB pipeline. The left column (Steps 1--7) processes forward vision and spatial scene geometry; the right column (Steps 8--14) executes cognitive driver assessment, thermal state integration, Pacejka braking dynamics, regulatory NCAP scoring, and Bird's-Eye View (BEV) radar synthesis.}
  \label{fig:pipeline_flowchart}
\end{figure*}

The proposed software framework processes continuous digital video streams through a modular, fourteen-step execution pipeline structured into two coupled columns: a \textbf{Perception and Spatial Scene Understanding} pipeline (Steps~1--7) and a \textbf{Physical Decision and Dynamic Actuation} pipeline (Steps~8--14), as illustrated in Fig.~\ref{fig:pipeline_flowchart}.

\subsection{Perception Pipeline (Steps 1--7)}
\begin{enumerate}
  \item \textbf{Frame Acquisition}: Captures sequential RGB driving imagery from a forward monocular camera ($\ge 1080\text{p}$, 30\,fps).
  \item \textbf{Contrast Enhancement}: Executes Contrast Limited Adaptive Histogram Equalization (CLAHE) on the luminance ($V$) channel in the HSV color space under detected low-light or nocturnal conditions ($I_{\text{mean}} < 80$).
  \item \textbf{Multi-Class Neural Inference}: Ingests the normalized image tensor into a fine-tuned YOLOv8n network, generating bounding boxes, class indices, and confidence scores across Pedestrians, Cars, Bicycles, and Motorcycles.
  \item \textbf{SORT Target Association}: Tracks bounding box centroids across temporal frames using a 7-state linear Kalman filter coupled with Hungarian IoU assignment, establishing persistent track identities.
  \item \textbf{Monocular Metric Depth \& BEV Projection}: Maps 2D image coordinates to 3D vehicle-relative coordinates $(X_{\text{lat}}, Z_{\text{long}})$ using camera intrinsic parameters, mounting height, and anthropometric reference heights.
  \item \textbf{Ego-Corridor Spatial Filtering}: Detects lane boundaries via Hough line segmentation and constructs a dynamic trapezoidal safety corridor, filtering out peripheral non-colliding objects.
  \item \textbf{Environmental Classification}: Evaluates scene weather and road lighting (clear, rain, fog, night), outputting an environmental friction degradation factor $\mu_{\text{env}}$.
\end{enumerate}

\subsection{Decision \& Control Pipeline (Steps 8--14)}
\begin{enumerate}
  \setcounter{enumi}{7}
  \item \textbf{In-Cabin Driver Alertness Monitoring}: Analyzes driver facial landmarks via dlib and \texttt{compute\_ear}, classifying attentiveness into Alert, Tired, Distracted, Drowsy, or Microsleep.
  \item \textbf{Brake Rotor Thermal State Update}: Tracks kinetic energy dissipation and convective Newton cooling, computing active rotor temperature $T_{\text{rotor}}$ and friction fade multiplier $\mu_{\text{therm}}$.
  \item \textbf{Adaptive TTC Threshold Computation}: Fuses environmental conditions, driver cognitive state, and thermal fade into a composite scalar trigger threshold $\text{TTC}_{\text{thresh}}$.
  \item \textbf{AEB Activation Decision}: Compares target obstacle TTC against $\text{TTC}_{\text{thresh}}$; if $\text{TTC} < \text{TTC}_{\text{thresh}}$, initiates emergency braking sequence.
  \item \textbf{Pacejka Longitudinal Braking Simulation}: Integrates vehicle equations of motion at a 1\,ms timestep with dynamic front/rear weight transfer and closed-loop ABS wheel slip cycling.
  \item \textbf{Biomechanical HIC \& Euro NCAP Evaluation}: Computes residual impact velocity, calculates Head Injury Criterion (HIC), assigns AIS trauma classification, and scores Euro NCAP compliance.
  \item \textbf{BEV Radar Rendering \& Black-Box Telemetry}: Synthesizes a real-time HUD Bird's-Eye View polar radar map and writes telemetry buffers to the black-box incident logger.
\end{enumerate}

Algorithm~\ref{alg:aeb_pipeline} defines the formal pseudocode governing per-frame pipeline execution.

\begin{algorithm}[!t]
\caption{End-to-End Per-Frame AEB Execution Pipeline}
\label{alg:aeb_pipeline}
\begin{algorithmic}[1]
\REQUIRE Forward camera frame $I_t$, vehicle forward speed $v_t$, previous rotor temperature $T_{t-1}$, cabin face image $F_t$, sampling time $\Delta t$.
\ENSURE Braking command $a_{\text{cmd}}$, driver warning $W_t$, incident log entry $\mathcal{L}_t$.
\STATE $I'_t \leftarrow \text{CLAHE}(I_t)$ \textbf{if} $\text{mean\_intensity}(I_t) < 80$ \textbf{else} $I_t$
\STATE $\mathcal{D}_t \leftarrow \text{YOLOv8n\_Inference}(I'_t, \text{conf}=0.25, \text{iou}=0.45)$
\STATE $\mathcal{T}_t \leftarrow \text{SORT\_Update}(\mathcal{D}_t)$
\FOR{each active track $k \in \mathcal{T}_t$}
  \STATE $d_k, x_{\text{lat},k} \leftarrow \text{PinholeDepth}(y_{\text{bbox},k}, h_{\text{bbox},k})$
  \STATE $\dot{d}_k \leftarrow \text{KalmanVelocity}(\mathcal{T}_t, k)$
  \STATE $\text{TTC}_k \leftarrow d_k / |\dot{d}_k|$ \textbf{if} $\dot{d}_k < 0$ \textbf{else} $\infty$
\ENDFOR
\STATE $\mathcal{C}_{\text{ego}} \leftarrow \text{ConstructEgoCorridor}(I'_t)$
\STATE $\mathcal{T}_{\text{in-path}} \leftarrow \{k \in \mathcal{T}_t \mid (x_{\text{lat},k}, d_k) \in \mathcal{C}_{\text{ego}}\}$
\STATE $S_{\text{driver}}, t_{\text{delay}} \leftarrow \text{AssessDriverFatigue}(F_t)$
\STATE $T_t \leftarrow \text{CoolDown}(T_{t-1}, v_t, \Delta t)$
\STATE $\mu_{\text{therm}} \leftarrow \text{FrictionMultiplier}(T_t)$
\STATE $\delta_{\text{thermal}} \leftarrow 0.5 \cdot (1.0 / \mu_{\text{therm}} - 1.0)$
\STATE $\text{TTC}_{\text{thresh}} \leftarrow 1.5 + \delta_{\text{weather}} + (t_{\text{delay}} - 0.7) + \delta_{\text{thermal}}$
\STATE $k^* \leftarrow \arg\min_{k \in \mathcal{T}_{\text{in-path}}} (\text{TTC}_k)$
\IF{$\text{TTC}_{k^*} < \text{TTC}_{\text{thresh}}$ \OR $S_{\text{driver}} = \text{MICROSLEEP}$}
  \STATE $a_{\text{cmd}} \leftarrow \text{ExecutePacejkaABS}(v_t, \mu_{\text{env}}, \mu_{\text{therm}})$
  \STATE $T_t \leftarrow T_t + \Delta E_{\text{kin}} / (m_{\text{rotor}} c_p)$
  \STATE $W_t \leftarrow \text{EMERGENCY\_BRAKE\_ACTIVE}$
\ELSIF{$\text{TTC}_{k^*} < \text{TTC}_{\text{thresh}} + 0.8$}
  \STATE $a_{\text{cmd}} \leftarrow 0.0$; $W_t \leftarrow \text{COLLISION\_WARNING}$
\ELSE
  \STATE $a_{\text{cmd}} \leftarrow 0.0$; $W_t \leftarrow \text{NORMAL\_DRIVING}$
\ENDIF
\STATE $\mathcal{L}_t \leftarrow \text{LogTelemetry}(\text{TTC}_{k^*}, v_t, T_t, S_{\text{drv}}, a_{\text{cmd}})$
\RETURN $a_{\text{cmd}}, W_t, \mathcal{L}_t$
\end{algorithmic}
\end{algorithm}

% ==============================================================================
% Section 4: Perception Layer
% ==============================================================================
\section{Perception, Scene Understanding, and Explainability}
\label{sec:perception}

\subsection{Dataset Curation and Neural Architecture}
The deep perception model is trained on a curated 6,000-image subset extracted from the Berkeley DeepDrive (BDD100K) dataset~\cite{yu2020}. BDD100K encompasses wide geographic, meteorological, and diurnal diversity across urban, residential, and highway environments. Ground-truth bounding-box annotations were extracted and mapped into four primary ADAS object classes: Pedestrian, Car, Bicycle, and Motorcycle.

The 6,000-image dataset was partitioned using stratified sampling into an 80\% training split (4,800 images) and a 20\% validation split (1,200 images). Severe class imbalance is an inherent characteristic of real-world driving scenes: passenger cars outnumber vulnerable road users by more than 3:1. To prevent gradient starvation across underrepresented VRU classes, training employed mosaic data augmentation, random affine scaling, horizontal flips ($p=0.5$), and automated mixup.

We fine-tuned the YOLOv8n (nano) architecture utilizing PyTorch with an SGD optimizer (initial learning rate $\text{lr}_0 = 0.01$, momentum $= 0.937$, weight decay $= 0.0005$). The model incorporates three decoupled detection heads predicting bounding-box offsets via Complete IoU (CIoU) loss, object classification via Binary Cross-Entropy (BCE) loss, and distribution focal loss (DFL) for box boundary refinement. Training executed over 30 complete epochs on an NVIDIA Tesla T4 GPU (total compute duration: 3,301.05 seconds).

\subsection{Empirical Detection Metrics and Failure Analysis}
In contrast to earlier unverified preliminary drafts that asserted fabricated mAP scores exceeding 77--89\%, Table~\ref{tab:verified_detection} reports the exact, verified numerical outputs extracted directly from the training logs (\texttt{results/training\_multiclass/results.csv}, Epoch 30).

\begin{table}[!t]
\caption{Verified YOLOv8n Multi-Class Detection Performance on BDD100K Validation Split (Epoch 30 Final)}
\label{tab:verified_detection}
\centering
\small
\resizebox{\columnwidth}{!}{%
\begin{tabular}{lcccc}
\toprule
\textbf{Metric} & \textbf{Verified Value} & \textbf{Original Claim} & \textbf{Gap} \\
\midrule
Overall mAP@0.5      & \textbf{40.01\%} (0.40008) & 77.10\% & $-37.09\%$ \\
Overall mAP@0.5:0.95 & \textbf{21.28\%} (0.21279) & 44.80\% & $-23.52\%$ \\
Overall Precision    & \textbf{53.60\%} (0.53601) & 78.50\% & $-24.90\%$ \\
Overall Recall       & \textbf{37.88\%} (0.37879) & 75.90\% & $-38.02\%$ \\
\midrule
Validation Box Loss  & 1.3055 & --- & --- \\
Validation Cls Loss  & 0.7142 & --- & --- \\
Validation DFL Loss  & 0.9879 & --- & --- \\
\bottomrule
\end{tabular}
}%
\end{table}

The empirical overall mAP@0.5 of 40.0\% reflects the physical reality of resource-constrained embedded edge backbones. As documented in the repository's baseline evaluation (\texttt{README.md}), class-specific mAP50 benchmarks are approximately 0.68 for Cars, 0.43 for Pedestrians, 0.23 for Bicycles, and 0.27 for Motorcycles. When trained exclusively as a single-class pedestrian detector on 3,400 images (\texttt{notebooks/yolo\_training.ipynb}, Cell 24), the model achieves an mAP@0.5 of 61.4\%, precision of 73.1\%, and recall of 52.5\%. The degradation observed in multi-class deployment stems from:
\begin{enumerate}
  \item \textbf{Backbone Capacity Constraint}: YOLOv8n possesses only 3.2\,million parameters. Resolving fine-grained structural differences between bicycles and motorcycles under small pixel footprints ($< 32\times32$ pixels at distances $> 25$\,m) saturates the feature representation capacity of the shallow C2f backbone.
  \item \textbf{Recall Deficit and Downstream Safety Implications}: The overall recall of 37.9\% indicates that approximately 62\% of distant or partially occluded target instances fail to register a detection above the confidence threshold ($\text{conf} \ge 0.25$). In an AEB safety context, a missed detection eliminates early TTC warnings. We analyze this vulnerability explicitly in Section~\ref{sec:limitations}, demonstrating how Kalman-filter state covariance extrapolation mitigates single-frame detection drops.
\end{enumerate}

Fig.~\ref{fig:pr_f1_curves} displays the verified Precision-Recall and F1-confidence curves across all four target classes.

\begin{figure}[!t]
  \centering
  \includegraphics[width=0.48\columnwidth]{training_04_PR_curve}\hfill
  \includegraphics[width=0.48\columnwidth]{training_05_F1_curve}
  \caption{Empirical verification curves extracted from Ultralytics training logs: (left) Precision-Recall profile across classes; (right) F1 score vs. confidence threshold, establishing an optimal operating point at confidence $\approx 0.28$.}
  \label{fig:pr_f1_curves}
\end{figure}

\subsection{Pinhole Geometry and Bird's-Eye View Mapping}
Longitudinal metric distance ($Z$) is reconstructed via the optical pinhole projection geometry:
\begin{equation}
  Z \;=\; \frac{f_y \cdot H_{\text{real}}}{h_{\text{bbox}}}, \qquad
  X_{\text{lat}} \;=\; \frac{(u_c - c_x) \cdot Z}{f_x},
  \label{eq:pinhole}
\end{equation}
where $f_x, f_y$ represent camera focal lengths in pixels ($f_x = f_y \approx 718.8$\,px for typical automotive calibration), $(c_x, c_y)$ is the optical principal center, $H_{\text{real}}$ is the assumed standard vertical height of the detected object class ($H_{\text{ped}} = 1.75$\,m, $H_{\text{car}} = 1.50$\,m, $H_{\text{child}} = 1.15$\,m), and $h_{\text{bbox}}$ is the observed pixel height of the bounding box. Transformed coordinates are projected onto a top-down polar Bird's-Eye View (BEV) radar display, mapping target trajectories in Cartesian meters relative to the front bumper.

\subsection{Regulatory Traffic Sign Recognition Layer}
To ensure that emergency braking thresholds conform to municipal speed regulations, we integrated a dedicated traffic sign recognition module (\texttt{src/traffic\_signs/sign\_detector.py}). The system monitors roadside regulatory circles and speed limit markers, extracting active zone limits:
\begin{equation}
  \gamma_{\text{zone}} \;=\;
  \begin{cases}
    1.50, & v_{\text{limit}} \le 30\,\text{km/h}, \\
    1.00, & 30 < v_{\text{limit}} \le 50\,\text{km/h}, \\
    0.90, & v_{\text{limit}} \ge 80\,\text{km/h}.
  \end{cases}
  \label{eq:sign_scaling}
\end{equation}
In pedestrian-dense school zones ($30\,\text{km/h}$), the threshold multiplier $\gamma_{\text{zone}} = 1.50$ expands the nominal TTC trigger buffer by 50\%, forcing the system to intervene much earlier in high-risk zones.

\subsection{Perceptual Explainability via Grad-CAM}
Autonomous braking requires complete algorithmic accountability. Black-box deep neural networks risk making safety-critical decisions based on visual background artifacts (e.g., streetlamps or asphalt markings). To verify that our detector focuses on genuine anatomical VRU features, we implemented Gradient-weighted Class Activation Mapping (Grad-CAM)~\cite{selvaraju2017} on the terminal C2f feature maps (\texttt{src/explainability/gradcam.py}).

The importance weight $\alpha_k^c$ for convolutional feature map $A^k$ with respect to class score $y^c$ is evaluated via global average pooling of spatial gradients:
\begin{equation}
  \alpha_k^c \;=\; \frac{1}{Z_{\text{pixels}}} \sum_{i} \sum_{j} \frac{\partial y^c}{\partial A_{i,j}^k}.
  \label{eq:gradcam_weights}
\end{equation}
The localized visual saliency heatmap $L_{\text{Grad-CAM}}^c$ is computed as:
\begin{equation}
  L_{\text{Grad-CAM}}^c \;=\; \text{ReLU}\left( \sum_k \alpha_k^c A^k \right).
  \label{eq:gradcam_map}
\end{equation}
Grad-CAM activations confirm that the network concentrates its attention on the torso and bipedal lower extremities of pedestrians, validating that AEB triggers are driven by genuine VRU morphology.

% ==============================================================================
% Section 5: Driver Monitoring
% ==============================================================================
\section{In-Cabin Driver Monitoring and Cognitive State Estimation}
\label{sec:driver_monitoring}

\subsection{Eye Aspect Ratio Formulation}
Driver fatigue detection is implemented in \texttt{src/driver\_monitoring/} utilizing the \texttt{dlib} 68-point facial landmark predictor~\cite{king2009}. Following the geometric derivation of Soukupov\'a and \v{C}ech~\cite{soukupova2016}, landmark coordinates for the left eye ($p_{36}$--$p_{41}$) and right eye ($p_{42}$--$p_{47}$) are tracked continuously. The Eye Aspect Ratio (EAR) quantifies palpebral aperture:
\begin{equation}
  \text{EAR} \;=\; \frac{\|p_2 - p_6\| + \|p_3 - p_5\|}{2\,\|p_1 - p_4\|}.
  \label{eq:ear}
\end{equation}
The scalar EAR is invariant to in-plane head rotation and camera-to-driver distance.

\subsection{State Transition Logic and Verified Reaction Penalties}
In contrast to discrepancies between preliminary text drafts, Table~\ref{tab:verified_driver} documents the exact, verified configuration constants extracted from \texttt{configs/part5\_thresholds.yaml} and \texttt{src/driver\_monitoring/alertness\_state.py}.

\begin{table}[!t]
\caption{Verified Driver State Classification and Reaction Time Penalties}
\label{tab:verified_driver}
\centering
\small
\resizebox{\columnwidth}{!}{%
\begin{tabular}{llcc}
\toprule
\textbf{State} & \textbf{Temporal Condition} & \textbf{Delay $t_r$ (s)} & \textbf{AEB Penalty} \\
\midrule
\textbf{ALERT}      & $\text{EAR} \ge 0.25$, normal gaze     & 0.70\,s & Base (0.00\,s) \\
\textbf{TIRED}      & Frequent blinks, $>5$ in 2\,s          & 1.10\,s & $+0.40$\,s \\
\textbf{DISTRACTED} & Yaw $>30^\circ$ or Pitch $>25^\circ$   & 1.30\,s & $+0.60$\,s \\
\textbf{DROWSY}     & Contiguous $\text{EAR} < 0.18$ for $>0.5$\,s & 1.60\,s & $+0.90$\,s \\
\textbf{MICROSLEEP} & Contiguous $\text{EAR} < 0.18$ for $>1.5$\,s & $\infty$ & Override ($\text{Instant}$) \\
\bottomrule
\end{tabular}
}%
\end{table}

When the driver is alert, baseline human reaction delay is established at $t_{r,\text{alert}} = 0.70$\,s, during which the human retains manual control. If prolonged ocular closure indicates drowsiness ($\text{EAR} < 0.18$ for $>0.5$\,s), reaction latency increases to $1.60$\,s. To prevent this cognitive latency from eroding stopping distance, the vehicle advances its intervention point by adding a penalty $\delta_{\text{driver}} = t_r - 0.70$\,s to the AEB activation threshold. In the catastrophic event of microsleep ($>1.5$\,s), the system overrides throttle actuation and executes maximum emergency braking immediately.

Head pose orientation is simultaneously resolved by solving the Perspective-n-Point (solvePnP) problem across six key facial fiducials (nose tip, chin, eye corners, mouth corners), flagging distraction whenever yaw exceeds $\pm 30^\circ$ or pitch exceeds $\pm 25^\circ$.

% ==============================================================================
% Section 6: Vehicle Dynamics and Thermodynamics
% ==============================================================================
\section{Longitudinal Vehicle Dynamics and Brake Thermodynamics}
\label{sec:dynamics}

\subsection{Five-Force Longitudinal Equilibrium}
Longitudinal chassis deceleration is evaluated through dynamic force summation along the vehicle's body axis:
\begin{equation}
  M\,\dot{v} \;=\; -\left( F_{x,\text{f}} + F_{x,\text{r}} + F_{\text{aero}} + F_{\text{roll}} + F_{\text{grade}} \right),
  \label{eq:long_forces}
\end{equation}
where $M = 1,500\,\text{kg}$, and $F_{x,\text{f}}, F_{x,\text{r}}$ are front and rear tyre retarding forces. The non-tyre resistance terms are modeled physically as:
\begin{align}
  F_{\text{aero}} &\;=\; \frac{1}{2}\,\rho_{\text{air}}\,C_d\,A_f\,v^2, \label{eq:aero} \\
  F_{\text{roll}} &\;=\; C_{rr}\,M\,g\,\cos\theta, \label{eq:roll} \\
  F_{\text{grade}} &\;=\; M\,g\,\sin\theta, \label{eq:grade}
\end{align}
with air density $\rho_{\text{air}} = 1.225\,\text{kg/m}^3$, aerodynamic drag coefficient $C_d = 0.30$, frontal area $A_f = 2.20\,\text{m}^2$, rolling resistance $C_{rr} = 0.015$, and road grade angle $\theta$.

\subsection{Dynamic Axle Load Transfer}
Deceleration creates an inertial pitching moment that redistributes normal axle loads:
\begin{align}
  F_{z,\text{front}} &\;=\; M\,g\,\frac{b}{L} \;+\; \frac{h_{cg}}{L}\,M\,a_x, \label{eq:fz_front} \\
  F_{z,\text{rear}}  &\;=\; M\,g\,\frac{a}{L} \;-\; \frac{h_{cg}}{L}\,M\,a_x, \label{eq:fz_rear}
\end{align}
where wheelbase $L = 2.67$\,m, distance from center of gravity (CG) to front axle $a = 1.12$\,m, distance to rear axle $b = 1.55$\,m, and CG height $h_{cg} = 0.54$\,m (\texttt{configs/vehicle\_dynamics\_config.yaml}). Under maximum dry deceleration ($a_x \approx 0.85\,g$), front axle normal force increases from its static 58\% distribution to over 75\% of total vehicle weight, dramatically altering tyre force saturation.

\subsection{Pacejka '92 Tyre Model and ABS Modulation}
Longitudinal tyre traction coefficient $\mu_x(\kappa)$ is governed by the Pacejka '92 Magic Formula:
\begin{equation}
\begin{split}
  \mu_x(\kappa) \;=\; D\,\sin\Bigl( &C\,\arctan\bigl( B\kappa \\
  &- E(B\kappa - \arctan(B\kappa)) \bigr) \Bigr),
\end{split}
\label{eq:pacejka}
\end{equation}
where $\kappa = (v - \omega r) / v$ represents wheel longitudinal slip. Verified parameter sets for typical automotive passenger tyres across five road surfaces are defined in Table~\ref{tab:pacejka_params}.

\begin{table}[!t]
\caption{Verified Pacejka Magic Formula Parameters Across Road Surfaces}
\label{tab:pacejka_params}
\centering
\small
\resizebox{\columnwidth}{!}{%
\begin{tabular}{lccccc}
\toprule
\textbf{Surface} & \textbf{B (Stiffness)} & \textbf{C (Shape)} & \textbf{D (Peak)} & \textbf{E (Curvature)} & \textbf{Peak Slip $\kappa^*$} \\
\midrule
\textbf{Dry}    & 10.0 & 1.65 & 0.85 & 0.97 & 0.12 \\
\textbf{Wet}    & 12.0 & 1.50 & 0.55 & 0.95 & 0.11 \\
\textbf{Gravel} &  8.0 & 1.30 & 0.45 & 0.80 & 0.16 \\
\textbf{Snow}   &  6.0 & 1.20 & 0.25 & 0.70 & 0.20 \\
\textbf{Ice}    &  4.0 & 1.10 & 0.10 & 0.50 & 0.25 \\
\bottomrule
\end{tabular}
}%
\end{table}

\begin{figure}[!t]
  \centering
  \includegraphics[width=0.95\columnwidth]{pacejka_curves}
  \caption{Pacejka '92 Magic Formula friction-slip curves across five simulated road surfaces, highlighting peak adhesion regimes ($\kappa^* \approx 0.10$--$0.15$) and post-peak sliding degradation.}
  \label{fig:pacejka_curves}
\end{figure}

The integrated ABS controller monitors individual wheel slip $\kappa$. When slip exceeds the threshold $\kappa_{\text{thresh}} = 0.15$, hydraulic caliper pressure is released; when slip drops below the target $\kappa_{\text{target}} = 0.10$, pressure is reapplied at a 50\,ms cyclic frequency, preventing wheel lockup and maximizing tyre shear force.

\subsection{Brake Rotor Thermodynamics and Thermal Fade}
The thermal module (\texttt{src/vehicle\_dynamics/thermodynamics.py}) models brake rotor thermal dynamics under repeated braking cycles. Kinetic energy dissipated during deceleration is converted directly into rotor thermal energy:
\begin{equation}
  \Delta E_{\text{kin}} \;=\; \frac{1}{2}\,M\,\left( v_{\text{start}}^2 - v_{\text{end}}^2 \right).
  \label{eq:ke_diss}
\end{equation}
The bulk temperature spike across four ventilated cast-iron rotors ($m_{\text{rotor}} = 40\,\text{kg}$ total, specific heat $c_p = 450\,\text{J/(kg}\cdot\text{K)}$) is:
\begin{equation}
  \Delta T \;=\; \frac{\Delta E_{\text{kin}}}{m_{\text{rotor}}\,c_p}.
  \label{eq:delta_t}
\end{equation}
A single emergency stop from 80 to 0\,km/h dissipates $3.70 \times 10^5$\,J, raising the adiabatic bulk rotor temperature by $20.58^\circ\text{C}$. In aggressive driving scenarios with front brake bias (65\% force allocated to front rotors), local friction surface temperatures elevate at more than triple this rate.

Convective cooling during vehicle transit obeys Newton's cooling law:
\begin{equation}
  T(t + \Delta t) \;=\; T_{\text{env}} + \left( T(t) - T_{\text{env}} \right)\,\exp(-k_{\text{cool}}\,\Delta t),
  \label{eq:newton_cooling}
\end{equation}
where the cooling coefficient $k_{\text{cool}} = k_{\text{base}} + k_{\text{speed}}\,v$ ($k_{\text{base}} = 1.0 \times 10^{-4}\,\text{s}^{-1}$, $k_{\text{speed}} = 5.0 \times 10^{-5}\,\text{m}^{-1}$, $T_{\text{env}} = 20^\circ\text{C}$).

Friction fade multiplier $\mu_{\text{mult}}(T)$ follows empirical piecewise cast-iron thermal degradation (with temperature $T$ in $^\circ$C):
\begin{equation}
\begin{split}
  &\mu_{\text{mult}}(T) = \\
  &\begin{cases}
    1.00, & T < 300, \\
    1.00 - 0.002\,(T - 300), & 300 \le T < 450, \\
    0.70 - 0.002\,(T - 450), & 450 \le T < 650, \\
    \max(0.10,\, 0.30 - 0.001\Delta T), & T \ge 650,
  \end{cases}
\end{split}
\label{eq:fade_mult}
\end{equation}
where $\Delta T = T - 650$. To compensate for diminished braking force when rotors fade, the AEB controller applies an adaptive TTC buffer penalty for $\mu_{\text{mult}} < 0.95$:
\begin{equation}
  \delta_{\text{thermal}} \;=\; \min\left( 3.00,\, 0.50 \cdot \left( \frac{1}{\mu_{\text{mult}}} - 1.0 \right) \right).
  \label{eq:thermal_penalty}
\end{equation}
At $428^\circ\text{C}$ (simulated after six hard stops), friction degrades by 25.6\% ($\mu_{\text{mult}} = 0.744$), injecting a $+0.172$\,s penalty. At $650^\circ\text{C}$, friction drops by 70\% ($\mu_{\text{mult}} = 0.300$), enforcing a $+1.167$\,s threshold advance.

\begin{figure}[!t]
  \centering
  \includegraphics[width=0.95\columnwidth]{thermodynamic_fade}
  \caption{Simulated brake rotor thermal dynamics over repeated braking events, showing temperature spikes exceeding $450^\circ\text{C}$ and corresponding friction coefficient degradation.}
  \label{fig:thermo_fade}
\end{figure}

% ==============================================================================
% Section 7: Decision Layer & Recording
% ==============================================================================
\section{Intelligent Decision Layer, VRU Behavior, and Incident Recording}
\label{sec:decision}

\subsection{Composite Adaptive TTC Formulation}
Rather than relying on an invariant threshold, the AEB activation decision dynamically incorporates environmental, human, and thermal states:
\begin{equation}
  \text{TTC}_{\text{thresh}} \;=\; \text{TTC}_{\text{base}} + \delta_{\text{weather}} + \delta_{\text{driver}} + \delta_{\text{thermal}},
  \label{eq:composite_ttc}
\end{equation}
where $\text{TTC}_{\text{base}} = 1.50$\,s. Under ideal conditions (clear day, dry asphalt, alert driver, cold brakes), $\text{TTC}_{\text{thresh}} = 1.50$\,s. Under adverse conditions (fog/rain $\delta_{\text{weather}} = +0.50$\,s, drowsy driver $\delta_{\text{driver}} = +0.90$\,s, faded brakes $\delta_{\text{thermal}} = +1.17$\,s), the activation threshold expands to $4.07$\,s ($+171\%$), initiating vehicle deceleration well before the collision horizon closes.

\subsection{VRU Behavior Classification via Random Forest}
In safety-critical urban crossings, pedestrian motion intention fundamentally governs trajectory risk. We integrated a 5-class Random Forest behavior classifier (\texttt{src/intelligence/train\_behavior.py}, \texttt{models/behavior\_classifier.pkl}). 

The model was audited in Phase~2 (\texttt{verification/verify\_ml\_models.py}), verifying an ensemble of $N_{\text{trees}} = 200$ estimators trained across five ground-truth VRU kinematic states: Standing (0), Walking (1), Running (2), Crossing (3), and Approaching Road (4). The model utilizes 11 trajectory features extracted from the SORT tracking filter. Table~\ref{tab:rf_features} (left) documents the verified feature importances.

The empirical dominance of velocity features ($\bar{v}$, $v_{\text{max}}$, $v_x$) aligns with pedestrian crossing mechanics: pedestrians intending to cross display rapid acceleration perpendicular to the curb, differentiating them from stationary bystanders.

\subsection{Multi-Physics Adaptive AEB Decision Model and Hybrid Safety Arbitration}
While the deterministic composite threshold in Eq.~\eqref{eq:composite_ttc} establishes an analytic safety envelope, complex multi-dimensional edge cases---such as elevated rotor temperature coupled with partial driver distraction on intermediate-friction wet asphalt---benefit from continuous statistical risk arbitration. To address this, we developed and trained a multi-physics Adaptive AEB decision classifier (\texttt{src/intelligence/train\_adaptive\_aeb.py}, \texttt{models/adaptive\_aeb\_model.pkl}).

The classifier evaluates a six-dimensional multi-physics telemetry vector:
\begin{equation}
  x_{\text{aeb}} \;=\; [\text{TTC},\, d,\, v_{\text{rel}},\, \mu_{\text{env}},\, \text{EAR},\, T_{\text{rotor}}]^T,
  \label{eq:aeb_feature_vector}
\end{equation}
where $\text{TTC} \in [0.3, 6.5]$\,s, obstacle distance $d \in [1.0, 100.0]$\,m, relative approach speed $v_{\text{rel}} \in [2.0, 30.0]$\,m/s, road friction coefficient $\mu_{\text{weather}} \in [0.10, 0.85]$, driver Eye Aspect Ratio $\text{EAR} \in [0.08, 0.38]$, and rotor disc temperature $T_{\text{rotor}} \in [20, 720]^\circ$C.

The training dataset comprises $N = 12{,}000$ synthetic driving scenarios systematically sampled across operational design corridors. Ground-truth binary intervention labels ($y \in \{0, 1\}$, where 1 triggers emergency braking) were derived from the composite physical boundary defined in Eq.~\eqref{eq:composite_ttc}, strictly adhering to ISO~22839 forward collision mitigation standards. We trained a Random Forest ensemble of $N_{\text{trees}} = 250$ estimators with maximum tree depth constrained to 8 and balanced class weighting.

\begin{table}[!t]
\caption{Verified Gini Feature Importances of the Machine Learning Classifiers}
\label{tab:rf_features}
\centering
\small
\resizebox{\columnwidth}{!}{%
\begin{tabular}{llc|lc}
\toprule
\multicolumn{3}{c|}{\textbf{VRU Behavior Classifier (200 Trees)}} & \multicolumn{2}{c}{\textbf{Adaptive AEB Classifier (250 Trees)}} \\
\textbf{\#} & \textbf{Kinematic Feature} & \textbf{Gini} & \textbf{Multi-Physics Feature} & \textbf{Gini} \\
\midrule
1  & Mean velocity ($\bar{v}$)          & 0.2910 & Time-to-Collision ($\text{TTC}$)        & 0.6703 \\
2  & Ego distance ($d_{\text{ego}}$)   & 0.1457 & Distance to obstacle ($d$)              & 0.2051 \\
3  & Max velocity ($v_{\text{max}}$)   & 0.1295 & Driver Eye Aspect Ratio ($\text{EAR}$)  & 0.0503 \\
4  & Lateral velocity ($v_x$)          & 0.1053 & Relative approach speed ($v_{\text{rel}}$) & 0.0396 \\
5  & Longitudinal velocity ($v_y$)     & 0.0746 & Road weather friction ($\mu_{\text{env}}$) & 0.0195 \\
6  & Horizontal centroid ($c_x$)       & 0.0628 & Rotor temperature ($T_{\text{rotor}}$)   & 0.0151 \\
7  & Velocity std ($\sigma_v$)         & 0.0559 & ---                                     & ---    \\
8  & Lateral displacement ($\Delta x$)  & 0.0459 & ---                                     & ---    \\
9  & Net displacement ($\Delta s$)     & 0.0408 & ---                                     & ---    \\
10 & Longitudinal disp ($\Delta y$)    & 0.0374 & ---                                     & ---    \\
11 & Vertical centroid ($c_y$)         & 0.0111 & ---                                     & ---    \\
\bottomrule
\end{tabular}
}%
\end{table}

The model was audited in Phase~2 (\texttt{verification/verify\_ml\_models.py}, \texttt{verification/ml\_models\_verified.json}). As reported in Table~\ref{tab:rf_features} (right), the primary decision driver is $\text{TTC}$ (Gini importance $0.6703$), followed by longitudinal distance $d$ ($0.2051$) and driver cognitive state ($0.0503$). On the independent test split ($N_{\text{test}} = 1{,}800$ samples), the classifier achieved a 98.67\% validation agreement with ISO~22839 corridors, a test accuracy of 97.44\%, and balanced precision, recall, and F1-scores of 95.56\% (confusion matrix: 1,259 true negatives, 495 true positives, 23 false positives, and 23 false negatives).

To prevent single-point machine learning failures in safety-critical vehicle control, the controller implements a dual-channel hybrid arbitration strategy (\texttt{src/collision/aeb\_controller.py}):
\begin{align}
  \text{AEB}_{\text{cmd}} \;=\; &\text{AEB}_{\text{physics}} \;\lor \nonumber \\
  &\Big( \text{AEB}_{\text{ML}} \land [\text{TTC} \le 1.25\,\text{TTC}_{\text{thresh}}] \Big),
  \label{eq:hybrid_arbitration}
\end{align}
wherein the deterministic physical safety boundary ($\text{AEB}_{\text{physics}}$) serves as an inviolable fail-safe guarantee that triggers emergency deceleration regardless of ML inference, while the ML model proactively commands pre-braking when compound multi-physics indicators anticipate an impending collision.

\subsection{Trajectory Prediction Evaluation Metrics}
To evaluate obstacle horizon forecasting accuracy, the framework incorporates formal trajectory error metrics (\texttt{src/evaluation/metrics.py}):
\begin{align}
  \text{ADE} &\;=\; \frac{1}{N\,T} \sum_{i=1}^N \sum_{t=1}^T \|Y_{i,t} - \hat{Y}_{i,t}\|_2, \label{eq:ade} \\
  \text{FDE} &\;=\; \frac{1}{N} \sum_{i=1}^N \|Y_{i,T} - \hat{Y}_{i,T}\|_2, \label{eq:fde} \\
  \text{RMSE} &\;=\; \sqrt{\frac{1}{N\,T} \sum_{i=1}^N \sum_{t=1}^T \|Y_{i,t} - \hat{Y}_{i,t}\|_2^2}, \label{eq:rmse}
\end{align}
where $Y_{i,t}$ and $\hat{Y}_{i,t}$ denote ground-truth and predicted 2D coordinates across prediction horizon $T$.

\subsection{Black-Box Incident Recorder}
For regulatory compliance and post-crash accident reconstruction, the system implements a continuous Black-Box Incident Recorder (\texttt{src/recorder/black\_box.py}). The recorder maintains a circular ring buffer storing the most recent 300 video frames ($\approx 10$\,s at 30\,fps) along with timestamped vehicle telemetry (velocity, brake pressure, TTC, driver state, and rotor temperature). Upon AEB triggering or when minimum TTC falls below $1.5$\,s, the system automatically writes an incident package to disk, serializing pre-event and post-event trajectory logs in standardized JSON format alongside synchronized video clips.

% ==============================================================================
% Section 8: Results & Verification
% ==============================================================================
\section{Experimental Results and Physical Verification}
\label{sec:results}

\subsection{Longitudinal Braking Simulation vs. Kinematic Model}
Table~\ref{tab:stopping_distances} presents the verified stopping distance comparison across road surfaces, benchmarked against the simple kinematic model ($d = v^2/2\mu g$) and ISO~21994 reference corridors at $50\,\text{km/h}$ (see \texttt{stopping\_distance\_}\allowbreak\texttt{verified.csv}).

\begin{table*}[!t]
\caption{Verified Stopping Distance Benchmarks: Kinematic Formulation vs. Pacejka '92 + ABS Dynamic Simulation}
\label{tab:stopping_distances}
\centering
\small
\resizebox{\textwidth}{!}{%
\begin{tabular}{lccccccc}
\toprule
\textbf{Road Surface} & \textbf{Friction $\mu$} & \textbf{Speed} & \textbf{Kinematic ($d_{\text{kin}}$)} & \textbf{Pacejka Pure ($d_{\text{pure}}$)} & \textbf{With 0.33s Lag ($d_{\text{total}}$)} & \textbf{Kinematic Gap} & \textbf{ISO 21994 Corridor} \\
\midrule
\multirow{2}{*}{\textbf{Dry Asphalt}} & \multirow{2}{*}{0.85} & 50\,km/h  & 11.59\,m & 11.54\,m & \textbf{16.12\,m} & $-28.1\%$ & 13--22\,m \checkmark \\
                                      &                       & 100\,km/h & 46.34\,m & 45.75\,m & \textbf{54.92\,m} & $-15.6\%$ & 45--75\,m \checkmark \\
\midrule
\multirow{2}{*}{\textbf{Wet Asphalt}} & \multirow{2}{*}{0.55} & 50\,km/h  & 17.91\,m & 17.97\,m & \textbf{22.56\,m} & $-20.6\%$ & 22--35\,m \checkmark \\
                                      &                       & 100\,km/h & 71.62\,m & 70.86\,m & \textbf{80.03\,m} & $-10.5\%$ & 70--105\,m \checkmark \\
\midrule
\multirow{2}{*}{\textbf{Gravel}}      & \multirow{2}{*}{0.45} & 50\,km/h  & 21.89\,m & 27.22\,m & \textbf{31.81\,m} & $-31.2\%$ & 28--40\,m \checkmark \\
                                      &                       & 100\,km/h & 87.54\,m & 106.57\,m & \textbf{115.74\,m} & $-24.4\%$ & --- \\
\midrule
\multirow{2}{*}{\textbf{Snow}}        & \multirow{2}{*}{0.25} & 50\,km/h  & 39.39\,m & 37.74\,m & \textbf{42.33\,m} & $-7.0\%$ & 55--75\,m (Packed) \\
                                      &                       & 100\,km/h & 157.57\,m & 146.56\,m & \textbf{155.73\,m} & $+1.2\%$ & --- \\
\midrule
\multirow{2}{*}{\textbf{Ice}}         & \multirow{2}{*}{0.10} & 50\,km/h  & 98.48\,m & 115.07\,m & \textbf{119.65\,m} & $-17.7\%$ & 60--150\,m \checkmark \\
                                      &                       & 100\,km/h & 393.93\,m & 367.36\,m & \textbf{376.53\,m} & $+4.6\%$ & --- \\
\bottomrule
\end{tabular}
}%
\end{table*}

The simulation confirms that on dry asphalt at 50\,km/h, pure mechanical braking requires 11.54\,m. Accounting for 0.33\,s of sensor-to-actuator latency (camera exposure, neural inference, hydraulic valve rise time), the actual stopping distance expands to 16.12\,m. The simple kinematic model predicts only 11.59\,m, underestimating the necessary stopping corridor by 28.1\%. On gravel and wet asphalt, where tyre slip oscillations in the ABS loop dissipate retarding energy, the underestimation reaches 31.2\%. All simulated stopping distances fall strictly within ISO~21994 physical corridors.

\subsection{Euro NCAP Regulatory Protocol Verification}
We executed the full Euro NCAP Car-to-Pedestrian test suite across the four standardized scenarios:
\begin{enumerate}
  \item \textbf{CPFA-50}: Car-to-Pedestrian Farside Adult (50\% vehicle lateral overlap, walking at 5\,km/h).
  \item \textbf{CPNA-25}: Car-to-Pedestrian Nearside Adult (25\% vehicle lateral overlap, walking at 5\,km/h).
  \item \textbf{CPNC-50}: Car-to-Pedestrian Nearside Child emerging from behind an obscuring parked car (50\% overlap, running at 5\,km/h).
  \item \textbf{CPLA}: Car-to-Pedestrian Longitudinal Adult walking in the vehicle's travel lane.
\end{enumerate}

Each scenario was simulated across nine approach velocities: 20, 25, 30, 35, 40, 45, 50, 55, and 60\,km/h. Table~\ref{tab:ncap_summary} summarizes the verified results on dry and wet surfaces (\texttt{ncap\_matrix\_dry.csv} and \texttt{ncap\_matrix\_wet.csv}).

\begin{table}[!t]
\caption{Euro NCAP Test Matrix Outcome Summary Across 36 Simulated Cases}
\label{tab:ncap_summary}
\centering
\small
\resizebox{\columnwidth}{!}{%
\begin{tabular}{lcccc}
\toprule
\textbf{Surface Condition} & \textbf{Avoided ($\checkmark$)} & \textbf{Partial ($\odot$)} & \textbf{Collision ($\times$)} & \textbf{Overall Score} \\
\midrule
\textbf{Dry Asphalt} ($\mu=0.85$) & 34 (94.4\%) & 1 (2.8\%) & 1 (2.8\%) & \textbf{97.0\%} \\
\textbf{Wet Asphalt} ($\mu=0.55$) & 32 (88.9\%) & 1 (2.8\%) & 3 (8.3\%) & \textbf{92.0\%} \\
\bottomrule
\end{tabular}
}%
\end{table}

\begin{figure*}[!t]
  \centering
  \includegraphics[width=0.45\textwidth]{ncap_matrix_dry}\hfill
  \includegraphics[width=0.45\textwidth]{ncap_matrix_wet}
  \caption{Euro NCAP standardized Car-to-Pedestrian test matrix heatmaps across four scenarios and nine velocities: (left) Dry road conditions ($\mu=0.85$, 97\% composite score); (right) Wet road conditions ($\mu=0.55$, 92\% composite score). Green indicates complete collision avoidance; yellow denotes partial speed reduction; red indicates unmitigated impact.}
  \label{fig:ncap_matrices}
\end{figure*}

On dry asphalt, all scenarios achieve 100\% avoidance from 20 to 50\,km/h. The binding constraint is scenario \textbf{CPNC-50} (Child emerging from behind parked cars). Because the visual line-of-sight is physically occluded until the child clears the parked vehicle profile ($d_{\text{detect}} \approx 16.2$\,m), available braking distance is truncated. At 55\,km/h, AEB reduces impact speed from 55 to 21.4\,km/h (partial avoidance, 61\% score); at 60\,km/h, residual impact occurs at 32.8\,km/h. On wet surfaces ($\mu=0.55$), CPNC fails to avoid collision from 50\,km/h onward due to degraded tyre grip.

\subsection{Biomechanical Injury Mitigation and HIC Analysis}
Table~\ref{tab:hic_verified} documents the Head Injury Criterion (HIC) and Abbreviated Injury Scale (AIS) values evaluated across the velocity spectrum (see \texttt{hic\_biomechanics\_}\allowbreak\texttt{verified.csv}).

\begin{table}[!t]
\caption{Verified Head Injury Criterion (HIC) and Injury Severity vs. Impact Velocity}
\label{tab:hic_verified}
\centering
\small
\resizebox{\columnwidth}{!}{%
\begin{tabular}{ccccc}
\toprule
\textbf{Speed} & \textbf{HIC (Geometry)} & \textbf{Impact Zone} & \textbf{AIS Tier} & \textbf{Clinical Trauma Description} \\
\midrule
0\,km/h  &    0.0 & None       & 0 & No Impact \\
10\,km/h &   25.6 & Hood Mid   & 1 & Minor (Superficial contusions) \\
20\,km/h &  556.1 & Windshield & 3 & Serious (Concussion, non-fatal) \\
30\,km/h & 1532.4 & Windshield & 5 & Critical (Severe closed-head trauma) \\
40\,km/h & 3145.8 & Windshield & 6 & Fatal (Catastrophic cranial trauma) \\
50\,km/h & 5495.4 & Windshield & 6 & Fatal (Unsurvivable head acceleration) \\
\bottomrule
\end{tabular}
}%
\end{table}

\begin{figure}[!t]
  \centering
  \includegraphics[width=0.95\columnwidth]{hic_vs_speed}
  \caption{Analytical Head Injury Criterion (HIC) curve as a function of impact velocity, depicting transitions across AIS clinical trauma tiers from Minor (AIS~1) to Fatal (AIS~6).}
  \label{fig:hic_curve}
\end{figure}

The physical scaling of HIC obeys $a_{\text{peak}} \propto v$, which implies $\text{HIC} \propto v^{2.5}$. Under identical impact surface stiffness, reducing impact velocity from 50 to 30\,km/h produces a theoretical reduction factor of:
\begin{equation}
  \left(\frac{50}{30}\right)^{2.5} \;\approx\; 1.667^{2.5} \;\approx\; 3.59\times.
  \label{eq:hic_scaling}
\end{equation}
However, in vehicle-pedestrian collisions, vehicle front-end geometry introduces severe non-linearities: at 30\,km/h, pedestrian Wrap-Around Distance (WAD) deposits the head onto the compliant sheet-metal hood panel ($k \approx 50\,\text{kN/m}$), yielding low HIC values. At 50\,km/h, WAD projects the head directly into the laminated windshield glass and peripheral A-pillar frame ($k \approx 300\,\text{kN/m}$), causing HIC to spike dramatically from 1,532 to 5,495. Thus, even partial deceleration from 50 to 30\,km/h transforms an otherwise fatal accident into a survivable collision.

\subsection{Computational Runtime and Edge Latency Budget}
Table~\ref{tab:latency} details the per-frame computational latency measured on an NVIDIA Tesla T4 GPU (16\,GB VRAM) and Intel Xeon CPU (@ 2.20\,GHz).

\begin{table}[!t]
\caption{End-to-End Computational Latency Budget per 30\,fps Video Frame}
\label{tab:latency}
\centering
\small
\resizebox{\columnwidth}{!}{%
\begin{tabular}{lccc}
\toprule
\textbf{Pipeline Subsystem} & \textbf{Hardware} & \textbf{Latency (ms)} & \textbf{Budget Share} \\
\midrule
Frame Capture \& CLAHE        & CPU      & 3.2\,ms  & 6.1\% \\
YOLOv8n Neural Inference      & GPU (T4) & 38.5\,ms & 73.8\% \\
SORT Kalman Filter Tracking   & CPU      & 0.8\,ms  & 1.5\% \\
Pinhole Depth \& BEV Map      & CPU      & 0.4\,ms  & 0.8\% \\
Ego-Corridor \& Lane Filter    & CPU      & 2.1\,ms  & 4.0\% \\
Driver DMS (EAR + solvePnP)   & CPU      & 4.2\,ms  & 8.0\% \\
Brake Dynamics \& Thermo Loop & CPU      & 0.6\,ms  & 1.1\% \\
Decision \& Black-Box Logging  & CPU      & 2.4\,ms  & 4.6\% \\
\midrule
\textbf{Total End-to-End Latency} & \textbf{Hybrid} & \textbf{52.2\,ms} & \textbf{100.0\%} \\
\bottomrule
\end{tabular}
}%
\end{table}

The complete pipeline operates at approximately 19.2 frames per second (52.2\,ms total latency), fully satisfying real-time closed-loop safety constraints on embedded automotive computing platforms.

% ==============================================================================
% Section 9: Discussion & Limitations
% ==============================================================================
\section{Discussion, Limitations, and Threats to Validity}
\label{sec:limitations}

\subsection{Detection Uncertainty and Propagation into TTC Errors}
The honest evaluation of our perception model revealed an overall recall of 37.9\% across four classes, with lower performance on narrow-profile VRUs (Bicycles ~0.23 mAP50, Motorcycles ~0.27 mAP50). If a pedestrian remains undetected for multiple consecutive frames, the calculated TTC defaults to infinity ($\infty$). 

In our architecture, this failure mode is mitigated by SORT Kalman filtering: once a track identity is initialized, the constant-velocity state extrapolation sustains the predicted bounding-box position for up to $N_{\text{max\_age}} = 30$ frames (1.0\,s) during temporary visual occlusion or detector dropouts. However, initial track acquisition latency remains bounded by the detector's true positive rate. Scaling to larger parameter architectures (YOLOv8s or YOLOv8m) or multimodal fusion with automotive radar is essential for safety-critical deployment.

\subsection{Monocular Pinhole Scale Ambiguity}
The monocular depth estimation pipeline assumes a flat ground plane and standardized anthropometric object heights ($H_{\text{ped}} = 1.75$\,m). If a pedestrian is seated, kneeling, or of non-standard stature, longitudinal distance errors scale proportionally:
\begin{equation}
  \frac{\Delta Z}{Z} \;=\; \frac{\Delta H_{\text{real}}}{H_{\text{real}}}.
  \label{eq:depth_error}
\end{equation}
For a child of height 1.15\,m erroneously classified as an adult (1.75\,m), distance is overestimated by 52\%, severely delaying AEB activation. Our integration of class-specific height presets partially resolves this ambiguity, but true scale invariance necessitates stereo vision or solid-state LiDAR.

\subsection{Thermal Model Assumptions}
The lumped-capacitance rotor model assumes uniform bulk temperature throughout the disc. In physical emergency braking, rapid energy deposition generates steep cross-sectional temperature gradients, producing localized surface thermal hotspots that exceed bulk predictions by over 100$^\circ$C. Furthermore, pad friction recovery under cooling involves complex chemical re-curing that our piecewise empirical multiplier simplifies.

% ==============================================================================
% Section 10: Conclusion
% ==============================================================================
\section{Conclusion and Future Work}
\label{sec:conclusion}

This research developed, verified, and benchmarked an end-to-end physics-aware automatic emergency braking system for vulnerable road user protection. By directly coupling real-time computer vision with non-linear Pacejka tyre dynamics, brake rotor thermodynamics, in-cabin cognitive driver monitoring, and biomechanical injury modeling, the pipeline addresses the fundamental limitations of classical kinematic ADAS architectures.

Empirical verification established that:
\begin{enumerate}
  \item Fine-tuned YOLOv8n object detection achieves 40.0\% mAP@0.5 and 21.3\% mAP@0.5:0.95 across 6,000 BDD100K images, delivering 19.2\,fps edge throughput on an NVIDIA Tesla T4.
  \item Classical kinematic models underestimate stopping distances by 15\% to 30\% by ignoring tyre slip saturation, ABS cycling, and dynamic load transfer.
  \item High-speed brake thermal fade elevates rotor temperatures to $428^\circ\text{C}$ after six stops, reducing friction by 25.6\% and mandating up to $+3.0$\,s in predictive TTC buffer penalties.
  \item Driver eye aspect ratio and head pose tracking successfully detect fatigue and microsleep, scaling reaction delays from 0.7 to 1.6\,s.
  \item Standardized Euro NCAP simulation yields 97\% and 92\% compliance scores on dry and wet surfaces across 36 test scenarios.
  \item Pre-crash partial braking from 50 to 30\,km/h drastically attenuates Head Injury Criterion (HIC) trauma, shifting impact outcomes from fatal cranial trauma down to survivable injuries.
  \item Multi-physics Random Forest threat classification achieves 98.67\% validation agreement with ISO~22839 corridors and 97.44\% test accuracy, integrated with dual-channel hybrid arbitration to guarantee fail-safe intervention.
\end{enumerate}

Future research will pursue: (i)~tight sensor fusion incorporating 77\,GHz FMCW automotive radar to eliminate monocular depth scale ambiguity; (ii)~embedded TensorRT INT8 quantization for sub-20\,ms perception latency; and (iii)~hardware-in-the-loop (HIL) chassis dynamometer validation.

% ==============================================================================
% Declarations & Acknowledgments
% ==============================================================================
\section*{CRediT Authorship Contribution Statement}
\textbf{Atharv Priyadarshi}: Conceptualization, Methodology, Software, Formal Analysis, Writing -- Original Draft.
\textbf{Agrawal Chehek Gopal}: Software, Data Curation, Validation, Investigation.
\textbf{Debangan Sarkar}: Software, Visualization, Resources, Data Curation.
\textbf{Arnav Garg}: Validation, Formal Analysis, Investigation.
\textbf{Indra Vir Singh}: Supervision, Project Administration, Methodology, Writing -- Review \& Editing.

\section*{Declaration of Generative AI in Scientific Writing}
During the preparation of this work, the authors utilized generative AI and assisted computational tools strictly for language clarity, stylistic formatting, and LaTeX source conversion. All scientific ideas, mathematical models, experimental scripts, numerical verification logs, and engineering conclusions were formulated, executed, and validated directly by the authors. The authors maintain full responsibility for the content of the published work.

\section*{Declaration of Competing Interest}
The authors declare that they have no known competing financial interests or personal relationships that could have appeared to influence the work reported in this paper.

\section*{Data and Code Availability}
All source code, trained neural network checkpoints, physics simulation scripts, and verification logs are publicly available in the project GitHub repository: \url{https://github.com/chehekagrawal/ADAS-Pedestrian-AEB}.

% ==============================================================================
% References
% ==============================================================================
\bibliographystyle{elsarticle-num}
\bibliography{ADAS_AEB_References}

\end{document}
